% Authoritative exercise source. Preserve all derivations and diagrams.

\ExerciseStart{G1.03}{Invertire un moto e ritrovare la preparazione}{Guida esplicita}

\BookHeading{Problema, dati e obiettivo}

Un corpo libero parte da \(x_0=1\,\mathrm m\) con velocità \(v=2\,\mathrm{m/s}\). Dopo \(T=3\,\mathrm s\) si inverte esattamente la velocità. Determinare la traiettoria di ritorno e l’istante in cui il corpo ritrova la posizione iniziale.

\ExTheory{\ExRef{c1-loschmidt}{1.02}}{La dinamica libera è \(x=x_0+vt\). L’inversione cambia \(v\) in \(-v\) conservando la posizione raggiunta. Conviene usare \(\tau=t-T\) per il tempo trascorso dopo l’inversione.}

\ExPrediction{Il ritorno ripercorre le stesse posizioni in ordine inverso. Nel grafico posizione-tempo le due pendenze hanno segni opposti; nel piano posizione-velocità l’intervento cambia la seconda coordinata.}

\ExFigure{\begin{tikzpicture}[>=Stealth,exlabel]
\node[exbox] (n0) at (0.0,0.0) {$(1,2)$};\node[exbox] (n1) at (3.6,0.0) {$(7,2)$};\draw[exother] (n0)--(n1);\node[exbox] (n2) at (7.2,0.0) {$(7,-2)$};\draw[exother] (n1)--(n2);\node[exbox] (n3) at (0.0,-1.6) {$(1,-2)$};\draw[exother] (n2.south)--++(0,-.4)-|(n3.north);\node[exbox] (n4) at (3.6,-1.6) {$(1,2)$};\draw[exother] (n3)--(n4);
\end{tikzpicture}}{Punti nel piano posizione-velocità. I cambi di segno sono interventi istantanei; gli altri passaggi sono evoluzioni libere.}


\begin{ExTasks}

\item Calcolare la posizione al tempo di inversione.

\item Scrivere la legge usando \(\tau\).

\item Calcolare le posizioni dopo uno e due secondi di ritorno.

\item Risolvere l’equazione di ritorno a \(x_0\).

\item Tradurre il risultato nel tempo totale e distinguere posizione e stato completo.

\end{ExTasks}

\ExSolution

\ExStep{Arrivare al punto di inversione}

\ExEq{x(T)&=1+2\cdot3=7\,\mathrm m.}
La coppia prima dell’intervento è \((x,v)=(7,2)\); immediatamente dopo è \((7,-2)\).

\ExStep{Impostare il ritorno}

\ExEq{x_R(\tau)&=7+(-2)\tau=7-2\tau,\qquad\tau\geq0.}
Il nuovo tempo parte da zero al momento dell’intervento.

\ExStep{Seguire gli istanti intermedi}

\ExEq{x_R(1)&=7-2=5,\qquad x_R(2)=7-4=3.}
All’andata si avevano \(x(2)=5\) e \(x(1)=3\): l’ordine dei punti è invertito.

\ExStep{Trovare il ritorno}

\ExEq{7-2\tau&=1,\quad -2\tau=-6,\quad\tau=3\,\mathrm s,\\t&=T+\tau=3+3=6\,\mathrm s.}
La durata della fase inversa deve uguagliare quella diretta.

\ExStep{Confrontare gli stati}

A \(t=6\) la posizione è 1, mentre la velocità è \(-2\). Si è ritrovata la configurazione spaziale con velocità invertite. Una seconda inversione porta \((1,-2)\) in \((1,2)\), ripristinando anche lo stato iniziale completo.

\par\noindent\begin{minipage}{\linewidth}
\ExFigure{\begin{tikzpicture}[>=Stealth,exlabel]
\draw[exaxis] (0,0)--(6.7,0) node[right] {$t$ (s)};\draw[exaxis] (0,0)--(0,3.3) node[above] {$x$ (m)};\draw[Accent,thick] (0,.4)--(3,2.8);\draw[Secondary,thick] (3,2.8)--(6,.4);\draw[dashed] (3,0)--(3,2.8);\node[below] at (3,0) {$3$};\node[below] at (6,0) {$6$};\node[left] at (0,.4) {$1$};\node[above] at (3,2.8) {$7$};
\end{tikzpicture}}{Traiettoria nel piano tempo-posizione. La cuspide identifica l’inversione, non un salto della posizione.}


\ExInterpret{Questo esempio libero rende esplicita la reversibilità. Nei sistemi a molte particelle il ritorno richiede l’inversione di tutte le velocità pertinenti e una dinamica compatibile con tale simmetria.}

\ExLab{Nella sezione 1.02 scegliere inversione ideale e lasciare alla fase inversa lo stesso tempo della fase diretta. Confrontare il ritorno spaziale delle particelle. Il calcolo qui svolto è il caso libero a una dimensione; urti e pareti della simulazione rendono le traiettorie più articolate.}


\end{minipage}
